% common/boxes.tex — entornos del dominio "taller de ciberseguridad".
% Los cuatro bloques que se repiten en cada sesión del curso:
%   caso  →  fuentes  →  pregunta/respuesta  →  mcq
% Depende de: tcolorbox, enumitem, xurl, environ, xstring.

\RequirePackage{environ}
\RequirePackage{xstring}

% ===== caso: enunciado del caso de estudio =====
% \begin{caso}{Ashley Madison (2015)} … \end{caso}
\newtcolorbox{caso}[1]{%
  enhanced, breakable, sharp corners,
  colback=accent!3, colframe=accent,
  boxrule=0.6pt, left=8pt, right=8pt, top=6pt, bottom=6pt,
  fonttitle=\bfseries\small, coltitle=white,
  attach boxed title to top left={xshift=6pt,yshift=-2pt},
  boxed title style={colback=accent, sharp corners, boxrule=0pt,
    top=1.5pt,bottom=1.5pt,left=5pt,right=5pt},
  title={\cysec@lbl{Caso}{Case}: #1}%
}

% ===== fuentes: lista de referencias del enunciado =====
% \begin{fuentes} \fuente{https://…} … \end{fuentes}
\newcommand{\fuente}[1]{\item \url{#1}}
\newenvironment{fuentes}{%
  \par\smallskip
  {\footnotesize\bfseries\color{accent}\cysec@lbl{Fuentes}{Sources}}%
  \begin{itemize}[leftmargin=1.4em,itemsep=1pt,topsep=2pt,font=\footnotesize]%
  \footnotesize
}{%
  \end{itemize}\par\smallskip
}

% ===== ¿está vacío el cuerpo de un entorno? =====
% Es lo que decide entre pintar la caja y dejar el aviso de «pendiente», y se
% pregunta una vez por cada `respuesta` y cada `analisis` del documento.
%
% Estaba con \StrSubstitute de xstring sobre el cuerpo entero detokenizado.
% xstring recorre la cadena carácter a carácter reconstruyéndola, así que el
% coste crece con el cuadrado del tamaño del cuerpo: un análisis de un párrafo
% largo costaba ~1,5 s **él solo**, y las once cajas del taller 2 eran 17 s de
% los 27 que tardaba el documento entero. Con l3tl es una pasada, y encima no
% hace falta ni empezarla cuando el cuerpo es largo: un cuerpo en blanco solo
% puede tener espacios y \par, o sea que largo y en blanco no puede estar.
\ExplSyntaxOn
\tl_new:N \l__cysec_probe_tl
% #1 cuerpo · #2 qué hacer si está en blanco · #3 si no
% `+m` y no `m`: los tres argumentos llevan párrafos dentro —el cuerpo del
% entorno y las dos ramas, que lo reimprimen—, y un argumento corto se planta
% en el primer \par con «Paragraph ended before … was complete».
\NewDocumentCommand \cysecIfBlankBody { +m +m +m }
  {
    \tl_set:Nx \l__cysec_probe_tl { \tl_to_str:n {#1} }
    \int_compare:nNnTF { \tl_count:N \l__cysec_probe_tl } > { 400 }
      { #3 }
      {
        \tl_remove_all:Nn \l__cysec_probe_tl { ~ }
        \exp_args:NNx \tl_remove_all:Nn \l__cysec_probe_tl { \tl_to_str:n { \par } }
        \tl_if_empty:NTF \l__cysec_probe_tl { #2 } { #3 }
      }
  }
\ExplSyntaxOff

% ===== pregunta / respuesta =====
\newcounter{pregunta}
\renewcommand{\thepregunta}{Q\arabic{pregunta}}

% \pregunta{¿Qué tipos de datos fueron expuestos?}
\newcommand{\pregunta}[1]{%
  \par\medskip
  \refstepcounter{pregunta}%
  \noindent{\bfseries\color{accent}\thepregunta.}\ {\bfseries #1}\par
}

% \begin{respuesta} … \end{respuesta}
% Cuerpo vacío ⇒ emite \todoans (caja en borrador + warning en el .log),
% de modo que `make ws-01 FINAL=1` delate lo que falta antes de entregar.
\NewEnviron{respuesta}{%
  \cysecIfBlankBody{\BODY}{%
    \todoans
  }{%
    \par\smallskip
    \begin{tcolorbox}[enhanced, breakable, sharp corners, colback=white,
      colframe=white, boxrule=0pt, leftrule=2pt, colframe=accent!35,
      left=10pt, right=2pt, top=2pt, bottom=2pt,
      before upper={\setlength{\parskip}{0.5em}}]%  tcolorbox lo pone a 0
      \BODY
    \end{tcolorbox}%
    \par\smallskip
  }%
}

% ===== mcq: selección múltiple ("1 o más respuestas válidas") =====
% \begin{mcq}{enunciado}
%   \opcion{…}          opción no elegida
%   \opcion*{…}         opción elegida (✓; resaltada con la opción `answers`)
% \end{mcq}
\RequirePackage{pifont}                      % \ding{51} = ✓ (fiable bajo XeTeX)
\newcommand{\cysec@check}{\ding{51}}
\newcommand{\opcion}{\@ifstar\cysec@optsel\cysec@optplain}
\newcommand{\cysec@optplain}[1]{\item #1}
% \item con etiqueta explícita NO incrementa el contador: hay que hacerlo a mano.
\newcommand{\cysec@optsel}[1]{%
  \refstepcounter{enumi}%
  \item[{\color{okgreen}\cysec@check}\;\alph{enumi}.]%
  \ifcysec@answers\color{okgreen}\bfseries\fi #1%
}

\newenvironment{mcq}[1]{%
  \par\medskip
  \refstepcounter{pregunta}%
  \noindent{\bfseries\color{accent}\thepregunta.}\ {\bfseries #1}\par
  \begin{enumerate}[label=\alph*.,leftmargin=2.2em,itemsep=2pt,topsep=3pt]%
}{%
  \end{enumerate}\par\smallskip
}

% ===== captura: evidencia gráfica del laboratorio =====
% \captura{archivo.png}{Título descriptivo de la captura}
% \captura[<opciones de \includegraphics>]{archivo.png}{Título}
% Si existe pics/archivo.png lo incluye como figura con su caption; si no,
% dibuja una caja punteada con el nombre exacto esperado (guardar la captura
% con ese nombre y recompilar basta para que aparezca). La label de la figura
% es fig:<archivo-sin-extensión>, p.ej. \cref{fig:sqlmap-hashes}.
% Un archivo faltante deja «Captura pendiente» en el .log también en FINAL,
% misma mecánica delatora que las respuestas vacías.
%
% El argumento opcional son las opciones que van dentro de los corchetes de
% \includegraphics, y por defecto es justo el ancho de siempre, así que las dos
% mil \captura ya escritas siguen saliendo idénticas. Existe para el recorte no
% destructivo del editor visual, que lo escribe como
% \captura[trim={.1\width .0\height .2\width .05\height},clip,width=0.8\linewidth]{…}{…}:
% el archivo de pics/ nunca se toca, el recorte es texto que se puede quitar.
\newcommand{\captura@base}{}
\newcommand{\captura}[3][width=0.8\linewidth]{%
  \IfSubStr{#2}{.}{\StrBefore{#2}{.}[\captura@base]}{%
    \renewcommand{\captura@base}{#2}}%
  \begin{figure}[htbp]
    \centering
    \IfFileExists{pics/#2}{%
      \includegraphics[#1]{pics/#2}%
    }{%
      \PackageWarning{cysec}{Captura pendiente: pics/#2}%
      \begin{tcolorbox}[enhanced, sharp corners, colback=gray!4,
        colframe=gray!55!black, boxrule=0pt,
        borderline={0.8pt}{0pt}{gray!55!black,
          dash pattern=on 4pt off 3pt},
        width=0.8\linewidth, boxsep=10pt, halign=center]
        {\bfseries\footnotesize\color{gray!55!black}%
          \cysec@lbl{CAPTURA PENDIENTE}{PENDING SCREENSHOT}}\\[4pt]
        {\footnotesize\color{gray!45!black}%
          \cysec@lbl{guardar como}{save as} \texttt{pics/#2}}%
      \end{tcolorbox}%
    }
    \caption{#3}%
    \label{fig:\captura@base}%
  \end{figure}%
}

% ===== analisis: interpretación de una captura o resultado =====
% Hermano de `respuesta` para lo que el profesor realmente evalúa: el
% análisis posterior a la evidencia. Cuerpo vacío ⇒ \todoanalisis (caja en
% borrador + warning «Análisis pendiente» en el .log, también en FINAL).
\newcommand{\todoanalisis}[1][]{%
  \PackageWarning{cysec}{Análisis pendiente\ifstrempty{#1}{}{ (#1)} en la
    captura o resultado más reciente}%
  \iffinalmode\else
    \par\smallskip
    \begin{tcolorbox}[enhanced, sharp corners, colback=yellow!8,
      colframe=orange!80!black, boxrule=0.4pt, left=6pt, right=6pt,
      top=3pt, bottom=3pt]%
      \small\bfseries\color{orange!60!black}%
      \cysec@lbl{Pendiente por analizar}{Analysis pending}%
      \ifstrempty{#1}{}{\normalfont\ — #1}%
    \end{tcolorbox}%
    \par\smallskip
  \fi
}

\NewEnviron{analisis}{%
  \cysecIfBlankBody{\BODY}{%
    \todoanalisis
  }{%
    \par\smallskip
    \begin{tcolorbox}[enhanced, breakable, sharp corners,
      colback=accent!3, colframe=accent!60,
      boxrule=0.4pt, left=8pt, right=8pt, top=8pt, bottom=6pt,
      fonttitle=\bfseries\small, coltitle=white,
      attach boxed title to top left={xshift=6pt,yshift=-2pt},
      boxed title style={colback=accent!60, sharp corners, boxrule=0pt,
        top=1.5pt,bottom=1.5pt,left=5pt,right=5pt},
      title={\cysec@lbl{Análisis}{Analysis}},
      before upper={\setlength{\parskip}{0.5em}}]%
      \BODY
    \end{tcolorbox}%
    \par\smallskip
  }%
}
